%!TEX TS-program = XeLaTeX
%use the following command:
%all document files must be coded in UTF-8
\documentclass[portuguese]{textolivre}
% build HTML with: make4ht -e build.lua -c textolivre.cfg -x -u article "fn-in,svg,pic-align"

\journalname{Texto Livre}
\thevolume{19}
%\thenumber{1} % old template
\theyear{2026}
\receiveddate{\DTMdisplaydate{2025}{11}{30}{-1}} % YYYY MM DD
\accepteddate{\DTMdisplaydate{2026}{2}{4}{-1}}
\publisheddate{\DTMdisplaydate{2026}{7}{23}{-1}}
\corrauthor{Eliane Schlemmer}
\articledoi{10.1590/1983-3652.2026.63158}
%\articleid{NNNN} % if the article ID is not the last 5 numbers of its DOI, provide it using \articleid{} commmand 
% list of available sesscions in the journal: articles, dossier, reports, essays, reviews, interviews, editorial
\articlesessionname{dossier}
\runningauthor{Schlemmer} 


\sectioneditorname{Francis Arthuso Paiva~\orcid{0000-0002-9083-3342}}
\layouteditorname{Rayany Chaves~\orcid{0009-0007-0028-484X}}


\title{“No princípio era a conexão”: linguagens e coinvenção no Paradigma da Educação OnLIFE}
\othertitle{“In the beginning was the connection”: languages and co-invention in the OnLIFE Education Paradigm}
% if there is a third language title, add here:
%\othertitle{Artikelvorlage zur Einreichung beim Texto Livre Journal}

\author[1,2]{Eliane Schlemmer~\orcid{0000-0001-8264-3234}\thanks{Email: \href{mailto:elianes@unisinos.br}{elianes@unisinos.br}}}
\affil[1]{Universidade do Vale do Rio dos Sinos, Escola de Humanidades, Programa de Pós-Graduação em Educação e Programa de Pós-Graduação em Linguística Aplicada, São Leopoldo, RS, Brasil.}
\affil[2]{Universidade Aberta, Laboratório de Educação a Distância e eLearning, Lisboa, Portugal.}

\addbibresource{article.bib}
% use biber instead of bibtex
% $ biber article

% used to create dummy text for the template file
\definecolor{dark-gray}{gray}{0.35} % color used to display dummy texts
\usepackage{lipsum}
\SetLipsumParListSurrounders{\colorlet{oldcolor}{.}\color{dark-gray}}{\color{oldcolor}}

% used here only to provide the XeLaTeX and BibTeX logos
\usepackage{hologo}

% if you use multirows in a table, include the multirow package
\usepackage{multirow}

% provides sidewaysfigure environment
\usepackage{rotating}

% CUSTOM EPIGRAPH - BEGIN 
%%% https://tex.stackexchange.com/questions/193178/specific-epigraph-style
\usepackage{epigraph}
\renewcommand\textflush{flushright}
\makeatletter
\newlength\epitextskip
\pretocmd{\@epitext}{\em}{}{}
\apptocmd{\@epitext}{\em}{}{}
\patchcmd{\epigraph}{\@epitext{#1}\\}{\@epitext{#1}\\[\epitextskip]}{}{}
\makeatother
\setlength\epigraphrule{0pt}
\setlength\epitextskip{0.5ex}
\setlength\epigraphwidth{.7\textwidth}
% CUSTOM EPIGRAPH - END

% to use IPA symbols in unicode add
%\usepackage{fontspec}
%\newfontfamily\ipafont{CMU Serif}
%\newcommand{\ipa}[1]{{\ipafont #1}}
% and in the text you may use the \ipa{...} command passing the symbols in unicode

% LANGUAGE - BEGIN
% ARABIC
% for languages that use special fonts, you must provide the typeface that will be used
% \setotherlanguage{arabic}
% \newfontfamily\arabicfont[Script=Arabic]{Amiri}
% \newfontfamily\arabicfontsf[Script=Arabic]{Amiri}
% \newfontfamily\arabicfonttt[Script=Arabic]{Amiri}
%
% in the article, to add arabic text use: \textlang{arabic}{ ... }
%
% RUSSIAN
% for russian text we also need to define fonts with support for Cyrillic script
% \usepackage{fontspec}
% \setotherlanguage{russian}
% \newfontfamily\cyrillicfont{Times New Roman}
% \newfontfamily\cyrillicfontsf{Times New Roman}[Script=Cyrillic]
% \newfontfamily\cyrillicfonttt{Times New Roman}[Script=Cyrillic]
%
% in the text use \begin{russian} ... \end{russian}
% LANGUAGE - END

% EMOJIS - BEGIN
% to use emoticons in your manuscript
% https://stackoverflow.com/questions/190145/how-to-insert-emoticons-in-latex/57076064
% using font Symbola, which has full support
% the font may be downloaded at:
% https://dn-works.com/ufas/
% add to preamble:
% \newfontfamily\Symbola{Symbola}
% in the text use:
% {\Symbola }
% EMOJIS - END

% LABEL REFERENCE TO DESCRIPTIVE LIST - BEGIN
% reference itens in a descriptive list using their labels instead of numbers
% insert the code below in the preambule:
%\makeatletter
%\let\orgdescriptionlabel\descriptionlabel
%\renewcommand*{\descriptionlabel}[1]{%
%  \let\orglabel\label
%  \let\label\@gobble
%  \phantomsection
%  \edef\@currentlabel{#1\unskip}%
%  \let\label\orglabel
%  \orgdescriptionlabel{#1}%
%}
%\makeatother
%
% in your document, use as illustraded here:
%\begin{description}
%  \item[first\label{itm1}] this is only an example;
%  % ...  add more items
%\end{description}
% LABEL REFERENCE TO DESCRIPTIVE LIST - END


% add line numbers for submission
%\usepackage{lineno}
%\linenumbers

\begin{document}
\maketitle

\begin{polyabstract}
\begin{abstract}
O artigo, orientado pelo Paradigma da Educação OnLIFE, propõe uma reflexão sobre a emergência de práticas linguageiras cocriadas na conexão entre humanos e Inteligências Artificiais Generativas (IAGen). Aborda a problemática da necessária superação da visão instrumental da IAGen na educação e das concepções tradicionais de linguagem, insuficientes diante das práticas discursivas híbridas que emergem com as IAGen. O objetivo é compreender como a cointeligência (conexão entre inteligências humanas e artificiais) atua como força coinventiva na produção de sentidos, problematizando a autoria e reconfigurando a função do docente, sobretudo, da área de linguagens. De caráter teórico-analítico, o texto dialoga com as epistemologias reticulares e as pedagogias conectivas, compreendendo a linguagem como acontecimento em rede, performado por humanos, algoritmos, dados, IAGen e plataformas enquanto agentes discursivos transorgânicos. Os resultados indicam que: 1) a IAGen, enquanto parceira epistêmica, em cointeligência, pode promover a coinvenção de mundos discursivos; 2) as práticas linguageiras que emergem em cointeligência tensionam as noções de autoria, criatividade, intencionalidade e dialogismo; 3) a necessidade de deslocar o docente da posição transmissiva para a de cocurador e intervencionista em fluxos discursivos, ao coabitar ecossistemas linguísticos híbridos e inteligentes, o que exige multiletramentos (digital, algorítmico e em IA) e transliteracia inseparáveis da formação crítica-sensível. Conclui-se que reconhecer as IAGen como agentes discursivos emergentes abre caminhos para uma educação linguística-plural-regenerativa, capaz de enfrentar desafios como a desinformação, os vieses algorítmicos e a perda de agência discursiva, além de potencializar a criação de mundos inventivos mais democráticos e inclusivos.

\keywords{Educação \sep Linguagem \sep Cultura digital \sep Comunicação digital \sep Inteligência Artificial}
\end{abstract}

\begin{english}
\begin{abstract}
Guided by the OnLIFE Education Paradigm, this article proposes a reflection on the emergence of co-created language practices arising from the connection between humans and Generative Artificial Intelligences (GenAI). It addresses the need to overcome the instrumental view of AI in education and the limitations of traditional conceptions of language, which are insufficient in the face of the hybrid discursive practices emerging with GenAI. The aim is to understand how co-intelligence (the connection between human and artificial intelligences) operates as a co-inventive force in meaning-making, problematizing authorship and reconfiguring the role of teachers, especially in the field of language studies.The text, theoretical-analytical in nature, engages with reticular epistemologies and connective pedagogies, understanding language as a networked event performed by humans, algorithms, data, GenAI, and platforms as transorganic discursive agents. The results indicate that: (1) GenAI, as an epistemic partner operating through co-intelligence, can foster the co-invention of discursive worlds; (2) language practices that emerge through co-intelligence challenge notions of authorship, creativity, intentionality, and dialogism; (3) there is a need to shift teachers from a transmissive stance to one of co-curatorship and intervention in discursive flows as they inhabit hybrid and intelligent linguistic ecosystems—an endeavor that requires multiliteracies (digital, algorithmic, and AI literacies) inseparable from critical-sensitive formation. It concludes that recognizing GenAI as an emerging discursive agent opens pathways to a plural, regenerative linguistic education capable of addressing challenges such as misinformation, algorithmic bias, and the loss of discursive agency, while enhancing the creation of inventive worlds that are more democratic and inclusive.

\keywords{Education \sep Language \sep Digital culture \sep Digital communication \sep Artificial Intelligence}
\end{abstract}
\end{english}
% if there is another abstract, insert it here using the same scheme
\end{polyabstract}

\section{O contexto e a virada epistêmica da conexão}\label{sec-intro}
A expressão “no princípio era a conexão”, mais do que uma metáfora, manifesta uma reinterpretação crítica da fundação logocêntrica ocidental – “no princípio era o verbo” –. Nessa provocação, o foco do “logos” (racional, isolado, humano), que posiciona a linguagem como origem do mundo e da razão, é deslocado para a conexão – uma ideia conectiva, distribuída, rizomática e não hierárquica, uma vez que, desde a concepção, existimos enquanto seres de e em conexão.

Sugere ainda um novo mito fundador: conhecimento e mundo emergem na/da conexão, na/da rede e na/da ecologia, o que implica um novo paradigma – o Paradigma da Educação OnLIFE –\cite{schlemmer2020ecossistemas,moreira2020educacaodigital,schlemmer2020ampliando,schlemmer2020onlife,schlemmer2024inteligencias}, cuja conexão não é somente digital ou técnica, mas da vida que conecta humanos e não humanos. A conexão, então, passa a ser o ato fundante em que emergem saberes, afetos, subjetividades e discursos. Pensar “no princípio” como conexão é compreender que a linguagem, o pensamento e o conhecimento emergem nos/dos atos conectivos transorgânicos,) \cite{difelice2017netativismo} em redes que interligam humanos, máquinas, ambientes e afetos.

O mundo tem sido marcado por pandemias (a exemplo da pandemia de SARS-COV-19), emergências climáticas extremas e a presença ubíqua de tecnologias digitais: plataformas, \textit{big data}, sensores, metaversos e, especialmente, a Inteligência Artificial Generativa (IAGen), o que evidencia que vivemos realidades hiperconectadas que desafiam os fundamentos criados num mundo pré-digital. Isso exige a problematização de aspectos ontológicos, antropológicos, sociológicos, epistemológicos e pedagógicos, uma vez que a experiência humana passa a operar cada vez mais distanciada dos modos de cognição exclusivamente humanos.

Historicamente, o pensamento antropocêntrico e dualista que governa o mundo tem operado a partir da centralidade no humano e na separação entre este e a máquina \cite{serres1993filosofiamestica,latour2012reagregando}. No entanto, Garry Kasparov, um dos maiores enxadristas do mundo, que em 1997 enfrentou o supercomputador Deep Blue da IBM em uma partida de xadrez histórica, inspirado em um ditado russo, “Se não pode vencê-los, junte-se a eles”, concebeu, em 1998, a ideia do “Xadrez Avançado”, no qual jogava com um computador, combinando a intuição humana com o cálculo da máquina. Essa vivência de parceria entre a inteligência humana e a inteligência da máquina levou Kasparov a afirmar que “A combinação de humanos e máquinas não é o futuro — é o presente. E essa colaboração nos torna mais fortes”\footnote{Em comunicação oral da divulgação do livro \textit{Deep Thinking: Where Machine Intelligence Ends and Human Creativity Begins} (publicado em maio de 2017).}.

A IA, enquanto campo de estudo dedicado à criação de máquinas capazes de realizar tarefas que, quando realizadas por pessoas, requerem inteligência, emergiu na década de 1950, sendo o termo “Inteligência Artificial” sugerido em 1956 por John McCarthy\footnote{Durante a Conferência de Dartmouth, que reuniu um grupo distinto de pesquisadores, incluindo John McCarthy, Marvin Minsky, Nathaniel Rochester e Claude Shannon, todos com um interesse comum: investigar o potencial das máquinas para simular a inteligência humana.}. Desde então, vimos surgir os primeiros sistemas especialistas\footnote{Capazes de decidir a partir de regras pré-estabelecidas.} (década de 1960), as redes neurais (RN)\footnote{Um tipo de algoritmo de aprendizado de máquina, inspirado na estrutura e funcionamento do cérebro humano e que deu origem ao primeiro Sistema de Tutor Inteligente (STI) – o SCHOLAR.} (década de 1970), as aplicações em diferentes contextos, a exemplo do xadrez (nas décadas de 1980 e 1990), ainda que marcadas por significativas limitações. Na década de 2000, devido ao crescimento exponencial da capacidade de processamento dos computadores e da disponibilidade de grandes quantidades de dados para treinar algoritmos, a IA atingiu um novo patamar de desenvolvimento. Já no século XXI, a partir de 2012, a IA começa a se popularizar e a evoluir rapidamente, impulsionada por avanços em áreas como aprendizado de máquina, redes neurais profundas, processamento de linguagem natural, visão computacional e robótica. Entretanto, foi em 2020 que IA começou a “ganhar o mundo”, com o desenvolvimento da Inteligência Artificial Generativa (IAGen), especialmente o ChatGPT, da OpenAI. Na sequência, várias IAGen surgiram, entre elas a DALL.E, Midjourney, Gemini, Llama etc., colocando a IA no centro das discussões globais, dada sua capacidade de gerar automaticamente conteúdo em diversos formatos (texto, imagem, áudio, código). Com o desenvolvimento da IAGen, fica evidente que o humano não é mais o centro exclusivo do conhecimento e da linguagem, uma vez que as diversas IAGen cocriam conosco textos, imagens, sons, vídeos, músicas, arte, códigos, espaços em 3D, revisões de literatura, narrativas multimodais, dentre outros.

No campo da Educação, a emergência das IAs, especialmente generativas, provoca debates sobre ensino, aprendizagem, autoria, criatividade, desenvolvimento da consciência crítica e ética, entre outros, uma vez que não apenas processam linguagem, mas integram processos de coinvenção de sentidos, reconfigurando os modos como aprendemos, ensinamos e produzimos conhecimento. Portanto, o desafio não se reduz a compreender como a IA entra na escola/universidade, mas a compreender como a escola/universidade se conecta com essa tecnologia e repensa a própria ideia de conhecimento e as formas como ele é produzido, o que exige mudança na cultura escolar/universitária e não apenas a inserção da IA nos processos de ensino e de aprendizagem, o que reforça a forma de operar já existente. A IA nos desafia a adotar um novo paradigma em educação, o que implica uma abordagem epistemológica que supere a visão antropocêntrica e dualista que separa humano-técnica-natureza \cite{serres1994contratonatural,latour2012reagregando}, e, dessa forma, contribua para repensar os espaços, tempos e formas de operar da educação formal, nos seus mais variados níveis \cite{schlemmer2023protagonismo,schlemmer2024inteligencias}.

Nesse contexto, a Educação OnLIFE \cite{schlemmer2020ecossistemas},compreendida como um percurso de pesquisa-desenvolvimento-formação, vem sendo constituída no âmbito do Grupo Internacional de Pesquisa Educação Digital – GPe-dU (UNISINOS/CNPq) e consolidou-se como paradigma durante a pandemia da Covid-19, tendo como proposição estruturante essa virada paradigmática e epistêmica. Em lugar de compreender as tecnologias digitais e a IA meramente como ferramentas a serem aplicadas e usadas no ensino, defende-se um deslocamento em direção a uma abordagem conectiva, orientada pela perspectiva da coinvenção \cite{schlemmer2024inteligencias}.

Essa virada paradigmática e epistêmica busca superar modelos de educação focados na representação ou na ação exclusivamente humana, que veem a máquina como ferramenta/recurso/meio/apoio a serem aplicados/usados a serviço de metodologias e práticas que pouco conversam com os novos habitares do ensinar e do aprender. Esses novos habitares emergem numa ecologia híbrida, que conecta escola, universidade, cidade, biodiversidade, tecnologias, plataformas e redes digitais, sistemas de geolocalização e de sensorização e a IA.

Neste artigo, propomos uma reflexão orientada pelo Paradigma da Educação OnLIFE, inscrita nas epistemologias reticulares \cite{difelice2012redessociais} e nas pedagogias conectivas \cite{schlemmer2023protagonismo} propiciada pela emergência de práticas linguageiras cocriadas entre humanos e IAGen. Articulamos uma leitura da linguagem como ato coinventivo, produzido em rede entre inteligências humanas e não humanas, sustentando a emergência de ecossistemas híbridos e inteligentes que ampliam o campo da educação linguística. O objetivo é compreender como a cointeligência, resultante da conexão entre inteligências humanas e artificiais, atua como força coinventiva na produção de sentidos, problematizando a autoria e reconfigurando a função docente e a do profissional da linguagem. Assim, ao articular “linguagens e coinvenção”, o título indica que: 1) a linguagem não é expressão de um sujeito, mas ato compartilhado e emergente; 2) a criação de sentido (discursiva, educativa, cultural) é coinventada por humanos e não humanos em simpoiese\footnote{Simpoiese (\textit{sympoiesis}), em que \textit{sym}=com; \textit{poiesis} = fazer/criar, significa literalmente “fazer-com”. Para \textcite{haraway2016staying}, simpoiese é o princípio segundo o qual a vida, o conhecimento e o mundo são produzidos em redes de interdependência radical, nas quais nenhum agente – humano, técnico ou biológico – existe ou age de forma isolada. O termo é proposto como contraponto do termo autopoiese \cite{maturana1980autopoiesis}, associada a sistemas que se auto-organizam de forma fechada.} \cite{haraway2016staying}; e  3) a linguagem, nesse contexto, é multimodal, sensível, híbrida e interdependente.

Do ponto de vista teórico, compreendemos que as IAGen configuram novas formas de vida discursiva, ao desafiarem as categorias modernas de sujeito, autoria e originalidade. Do ponto de vista prático, apresentamos experiências educativas desenvolvidas no âmbito da Rede Internacional de Educação OnLIFE (RIEOnLIFE)\footnote{A descrição detalhada, bem como as análises de algumas das principais e mais atuais experiências educativas podem ser encontradas em \textcite{schell2024tomkat,schlemmer2024universidade,thomazzoni2024vivencia,schell2025narrativas,schell2025doctoral}.}, nas quais estudantes, docentes e inteligências artificiais cocriaram textos, narrativas e reflexões relacionadas à linguagem, à ética e à invenção. A partir disso, apontamos pistas para a construção de uma Educação Linguística OnLIFE.

O texto está organizado em nove seções. A introdução (seção \ref{sec-intro}) apresenta o contexto, a problematização e os objetivos. A seção \ref{sec-normas} aborda o Paradigma da Educação OnLIFE, que orienta a discussão sobre “IAEducação?”, situando a presença da IA na Educação a partir de três abordagens que estruturam a seção \ref{sec-conduta}. A seção \ref{sec-4} discute a IAGen a partir da ontologia da conexão e o deslocamento da linguagem. A seção \ref{sec-5} trata das IAGen como agentes na coinvenção de mundos discursivos. Nesse sentido, problematiza-se a autoria e a ética da cointeligência, compondo a seção \ref{sec-6}. A seção \ref{sec-7} propõe a Educação Linguística OnLIFE e discute as implicações pedagógicas e a formação docente. A seção \ref{sec-8} apresenta as Considerações Finais, com pistas para uma Educação Linguística OnLIFE. Por fim, são inclusas as referências e os agradecimentos.

\section{O paradigma da educação OnLIFE}\label{sec-normas}
A presença de tecnologias, plataformas e redes digitais, em especial da IA, na sociedade e na educação implica uma transformação cultural profunda. Como afirma \textcite{Fagundes2011CulturaDigital}, a tecnologia digital amplia a cognição; o problema é como a escola vai se apropriar dessa tecnologia e integrar-se à cultura digital, pois não basta disponibilizar computadores em rede e formar professores para o uso/aplicação dessas tecnologias. É essencial mudar a cultura da escola/universidade. A cultura da escola/universidade não é a cultura digital. Esta só se estabelece quando se transformam, na escola e na universidade, os espaços, os tempos, o conceito de sala de aula e o seu funcionamento.

\begin{quote}
    Não é inclusão digital na escola [/universidade] é a inclusão da escola [/universidade] na cultura digital "A dificuldade da escola [/universidade] é a que a escola [/universidade] não transforma. É a organização mais resistente. O poder da escola [/universidade] está em não mudar \cite{Fagundes2011CulturaDigital}\footnote{Em comunicação oral no TEDx UNISINOS.}.
\end{quote}


Isso evidencia a necessidade de superar o paradigma educacional vigente e aponta para a urgência da construção de um novo paradigma, denominado Paradigma de Educação OnLIFE. 


O termo “onlife”, a partir de uma revisão de literatura realizada por \textcite{schlemmer2024inteligencias}, aparece inicialmente em \textcite{rheingold1993virtualcommunity} ao discutir as comunidades virtuais e novas formas de convivência que emergem nas redes, antecipando a ideia de uma experiência em que o \textit{online} e o \textit{offline} se entrelaçam. Em \textcite{gutowitz1995alifeonlife} que amplia o conceito ao propor \textit{The Alife/Onlife Service}, vinculando-o à vida artificial (\textit{A-Life}) e a ecossistemas digitais, nos quais agentes humanos e não humanos coexistem em rede, conferindo ao termo uma dimensão tecno-biológica e sistêmica. \textcite{becker2000onlifetraveler} retoma o termo “onlife” no contexto do metaverso AlphaWorld ao tratar do \textit{onlife traveler}" o viajante digital que habita mundos virtuais e vivencia a hibridização entre identidade, corpo e avatar, introduzindo um sentido estético e experiencial do viver em rede. No entanto, o termo ganha destaque com \textcite{floridi2015onlife}, em \textit{The Onlife Manifesto: Being Human in a Hyperconnected Era}, que o redefine em chave filosófica e ontológica para compreender a condição humana na hiperconectividade, na qual se dissolvem fronteiras entre real e virtual, humano e tecnológico, constituindo a Sociedade Onlife. Assim, o termo “onlife” percorre um caminho que vai da observação da sociabilidade digital à elaboração ontológica e ética de um novo modo de ser e conhecer em rede.

Na Educação, o termo “onlife” é apropriado e teorizado na tríade Pesquisa-Desenvolvimento-Formação que caracteriza o GPe-dU UNISINOS/CNPq, a partir do coengendramento de sentidos epistemológicos presentes nas pesquisas anteriores do Grupo e em continuidade à elaboração do conceito de Educação Híbrida e Multimodal \cite{schlemmer2016agamificacao,schlemmer2016bgames,schlemmer2016chibridismo}.Esse movimento provocou a necessidade de ampliar a compreensão das Epistemologias Reticulares e Conectivas \cite{difelice2012redessociais}, da Teoria da Complexidade \cite{morin2011setesaberes}, e da Teoria da Cognição Inventiva \cite{kastrup2015cognicao}, incorporando noções como simbiota, aprendizagem enquanto mestiçagem e invenção \cite{serres1993filosofiamestica}, ato conectivo transorgânico e transubstanciação \cite{difelice2017netativismo} simpoiésis \cite{haraway2016staying} e habitar atópico \cite{difelice2009paisagensposurbanas}, em diálogo com a Sociedade Onlife \cite{floridi2015onlife}. Com isso, propõe o conceito de Educação OnLIFE, que se consolida, diante dos desafios impostos pela pandemia da Covid-19, como um paradigma \cite{moreira2020educacaodigital,schlemmer2020ampliando,schlemmer2020onlife}.

O Paradigma da Educação OnLIFE rompe com a visão antropocêntrica e dualista, compreendendo a inseparabilidade entre humano, natureza e técnica, além de entender que o social se estende também aos não humanos \cite{latour2012reagregando}. Trata-se de uma educação conectada (\textit{on}) na vida (\textit{LIFE}), alinhada aos novos habitares digitais e híbridos e às novas linguagens – técnico-computacionais, midiáticas, multimodais, algorítmicas, generativas e ecossistêmicas — que reconfiguram comunicação, conhecimento, ensino e a aprendizagem. O digital e o híbrido deixam de apenas representar o mundo e passam a cocriar “comundos”. É a partir da problematização do mundo e comundos que os processos de ensino e de aprendizagem emergem e desenvolvem-se em atos que conectam humanos e não humanos num percurso de coinvenção em rede \cite{schlemmer2024inteligencias}. Esse movimento reticular e coinventivo problematiza a própria ideia de conhecimento e como é produzido, desestabilizando a supremacia da inteligência humana. Ao propor uma virada epistêmica, questiona a pré-existência de dois elementos — sujeito e objeto — que orientados por uma teoria da ação, seja de um sobre o outro ou mesmo de interação, promovem a formação, a partir de uma visão antropocêntrica e dualista do mundo, indicando a necessidade de migrar para a formação em contextos ecologicamente conectados.

Dessa forma, o paradigma tem potência de provocar a transubstanciação e a metamorfose da educação – seus espaços, tempos e formas de operar – impactando metodologias, práticas, currículos e formas de governança. Isso implica uma transformação qualitativa da educação, em vez de apenas transpor modelos antigos para o digital ou híbrido. 

Recentemente, frente às emergências climáticas extremas, especialmente as vivenciadas no Rio Grande do Sul\footnote{Segundo dados do Governo do Estado do RS, no período de abril e maio de 2024, devido às chuvas intensas e inundações, 478 dos 497 municípios do estado foram oficialmente afetados, o que corresponde a 96,2\% do território do RS. Mais de 2,4 milhões de pessoas afetadas direta ou indiretamente pelas enchentes e deslizamentos, cerca de 81 mil pessoas desabrigadas ou desalojadas no pico da emergência e cerca de 200 óbitos confirmados, caracterizando o evento como o maior desastre climático da história do estado em número de vítimas humanas.},  novos referenciais foram incorporados ao Paradigma da Educação OnLIFE, como o Pensamento Ecológico \cite{morton2010ecologicalthought}, o Ser Ecológico \cite{morton2018beingecological} e o conceito de Fitópolis \cite{mancuso2024fitopolis}. Eles ampliam a dimensão ontológica e ética do paradigma, orientando para uma educação ecológica e regenerativa, em que o aprender é também um ato de coabitar, eco-habitar e coevoluir com o planeta.

Isso nos leva a compreender que o Paradigma da Educação OnLIFE exige pedagogIAs\footnote{O termo PedagogIAs, do qual destaca-se a IA é abordado por \textcite{schlemmer2023protagonismo} para destacar a Pedagogia em tempos de Inteligências Artificiais.} conectivas que operam em rede com algoritmos, \textit{Big Data}, sensores, robôs, cidade, biodiversidade, clima, vírus e diversas inteligências. Nesse contexto, evidencia-se a dimensão ecológica-conectiva do ensino e da aprendizagem, em contraste com pedagogias centradas no diálogo exclusivo entre humanos. O objetivo passa a ser a construção e a significação de um novo modo de habitar o ensino e a aprendizagem em um comundo que se expande para além dos espaços geográficos, numa ecologia conectada \cite{schlemmer2023protagonismo}.

Assim, o Paradigma da Educação OnLIFE constitui-se a partir da epistemologia reticular, conectiva e coinventiva, que emerge das confluências entre tecnológico, ecológico e cognitivo, propondo uma educação viva, transorgânica e ecossistêmica, em que ensinar e aprender são atos de coinvenção, coexistência e corresponsabilidade no viver, conviver e habitar o mundo hiperconectado. Nesse contexto, como pensar a IAEducação?

\section{IAEducação? Do intrumental à cointeligência}\label{sec-conduta}

A presença da Inteligência Artificial na Educação (IAE) tem sido marcada por abordagens distintas, sistematizadas por \textcite{schlemmer2024inteligencias} no \textit{Continuum} IAEducação? (\Cref{fig1}). Em especial, o avanço da IAGen, popularizada em 2022, exige o deslocamento de um enfoque instrumental para um conectivo – a CoInteligência – alinhado ao Paradigma da Educação OnLIFE.

\begin{figure}[h!]
\centering
\begin{minipage}{.9\textwidth}
 \includegraphics[width=\textwidth]{Fig1.png}
 \caption{\textit{Continuum} IAEducação?.}
 \label{fig1}
 \source{Elaboração própria a partir de comunicação oral \cite{schlemmer2024inteligencias}.}
\end{minipage}
\end{figure}
 
O “\textit{Continuum} IAEducação?” (\Cref{fig1}) apresenta três abordagens: 1) Abordagem Instrumental; 2) Abordagem Relacional; e 3) Abordagem da CoInteligência.

A Abordagem Instrumental compreende a IA como ferramenta a ser usada/aplicada na Educação. Essa abordagem entende a IA primariamente como uma ferramenta ou, em casos extremos, como um substituto do professor. O foco está na eficiência do ensino: analisar dados sobre o desempenho dos estudantes e identificar áreas nas quais precisam de apoio, personalizar o ensino, automatizar tarefas administrativas e criar sistemas de tutoria inteligentes que forneçam feedback personalizado e adaptam o conteúdo às necessidades do estudante. Baseia-se numa visão antropocêntrica e dualista, que compreende o conhecimento como externo ao sujeito (epistemologia empirista) e sustenta pedagogias diretivas. \textcite{schlemmer2024inteligencias} cita o \textit{David Game College}, no Reino Unido, que usa IA em substituição a professores humanos. Essa abordagem limita o potencial transformador da IA ao reforçar currículos instrucionistas e modelos centrados na transmissão e memorização de informações.

A Abordagem Relacional é considerada por \cite{schlemmer2024inteligencias} como intermediária, pois se apropria da IA na relação com a educação, defendendo o desenvolvimento de competências digitais e a Didática com IA, como propõem documentos elaborados pela UNESCO e pela União Europeia. Embora essa abordagem reconheça o conhecimento como interação e avance para uma epistemologia interacionista, ainda se mantém a separação entre inteligências humanas e artificiais, tratando a IA predominantemente como ferramenta e preservando traços de antropocentrismo e dualismo, expressos em pedagogias relacionais.

A Abordagem da Cointeligência – “Inteligências Conectadas”– é defendida e trabalhada na perspectiva do Paradigma da Educação OnLIFE \cite{schlemmer2024inteligencias}. Ao compreender a IA como parceira epistêmica – uma inteligência de natureza distinta que se articula à inteligência humana em cointeligência – torna-se possível promover a coinvenção de mundos discursivos, superando visões antropocêntricas e dualistas. A cointeligência refere-se a uma forma de inteligência coengendrada que, pelo ato conectivo entre humano e não humano (transorgânico), produz um conhecimento que é coinventado, resultando em um desempenho superior ao que seria alcançado individualmente. Essa abordagem é orientada pelas epistemologias reticulares \cite{difelice2012redessociais} e expressa em pedagogias conectivas \cite{schlemmer2023protagonismo} – no lugar de configurar uma pedagogia específica (Pedagogia da IA), bem como por metodologias e práticas inventivas \cite{schlemmer2022metodologias}. Nesse contexto, a IA deixa de ser pensada enquanto competência individual, separada das demais, e passa a integrar a Rede de Competências para a docência OnLIFE \cite{schlemmer2018projetos,schlemmer2022metodologias}. \textcite{mollick2024cointelligence} afirma que a produção intelectual cada vez mais será produto de cointeligência e não de comunidades. A abordagem da cointeligência representa uma oportunidade efetiva de superação do paradigma educacional vigente, ao instigar o estudante a deslocar o foco da resposta para a pergunta, para a problematização e para a coinvenção em rede.Trata-se de um movimento de agenciamento entre inteligência diversas, que pelo ato conectivo transorgânico \cite{difelice2017netativismo} produz uma ecologia de inteligências. Essa abordagem leva à emergência de hiperinteligências \cite{lovelock2019novacene}.


O conceito de HiperInteligêncIA, apresentado por \textcite{lovelock2019novacene} em \textit{Novaceno}, compreende a inteligência humana como parte e complemento de outras inteligências (das quais ela depende) e que irão sucedê-la. Assim, a hiperinteligência desafia a visão antropocêntrica e dualista que separa inteligência humana (cerebrocêntrica) da inteligência artificial (IA) e de outros tipos de inteligências não cerebrocêntricas, como a inteligência das plantas \cite{mancuso2017plantrevolution} e a inteligência dos fungos \cite{sheldrake2020entangledlife}, as quais são constituídas por arquiteturas reticulares e distribuídas e que implicam redes simpoiéticas. A emergência das hiperinteligências implica alianças entre humanos e não humanos, os quais se hibridizam em novas potências conectivas. Não se trata de uma simples soma de inteligências distintas, mas de uma ecologia conectiva que, sendo simpoiética, produz \textit{comundos}, hibridismos e invenções.


Dessa forma, podemos compreender que a abordagem da cointeligência implica: 1) Inteligências Conectadas: referem-se às diferentes inteligências (humanas e não humanas) em ato conectivo transorgânico; 2) CoInvenção: processo de invenção entre diferentes inteligências conectadas, que se produz enquanto invenção conectada e 3) HiperInteligêncIA: não é uma superinteligência, mas emerge na expansão da conectividade, a qual propicia a hibridização de diferentes inteligências conectadas em coinvenção, deslocando a discussão para uma perspectiva ecológica-conectiva.

Dessa maneira, a abordagem da cointeligência, orientada pelo Paradigma da Educação OnLIFE, problematiza a supremacia da inteligência humana e aponta para a necessidade de migrar da formação de sujeitos para a formação em contextos ecologicamente conectados, nos quais inteligências diversas operam em coinvenção, potencializando a emergência das hiperinteligências.

Assim, quando problematizamos a IAEducação?, não estamos falando da IA como ferramenta, recurso, meio, apoio a ser usado/aplicado na educação, tampouco de ensinar sobre a IA na Educação (usos e aplicações), ou ainda da IA e/na/para a Educação, numa perspectiva relacional; e sim da IA como parceira epistêmica, na qual a inteligência humana e a IA estão coengendradas.

\section {As IA Generativas, a ontologia da conexão e o deslocamento da linguagem}\label{sec-4}

As IAGen (como \textit{Large Language Models – LLMs}) funcionam como sequenciadores linguístico-probabilísticos de baixa entropia cruzada. São modelos matemáticos que predizem a sequência probabilística de palavras, sendo resultado de uma complexa orquestração de \textit{softwares}, dados, algoritmos de aprendizado profundo e \textit{feedback} humano. Nesse sentido, o desenvolvimento da IA Gen faz parte de uma complexidade que, além de ser conectiva, é simpoiética (\textit{fazer-junto}) \cite{haraway2016staying} e atópica \cite{difelice2012redessociais}. Há, portanto, uma implicação radical de que nada se faz por si mesmo, sendo a cocriação com dados e algoritmos essencial para a produção do mundo que habitamos.

\textcite{accoto2020mundodado} afirma que “o mundo dado é o mundo dos dados”, deslocamento que rompe com a concepção representacional da técnica e nos provoca a uma ontologia em que o digital não espelha o real, mas o coconstrói. Nesse entendimento, dados, algoritmos e infraestruturas digitais deixam de ser descrições ou mediações neutras e passam a operar como agentes coformadores do mundo, constituindo práticas, percepções e modos de vida. Dessa forma, a linguagem não pode mais ser compreendida como instrumento de representação, mas como acontecimento – um fluxo inventivo que se atualiza na (e pela) própria rede, emergindo no entrelaçamento entre inteligências humanas, artificiais, ambientais e técnicas. As palavras, as imagens e os códigos deixam de ocupar esferas distintas e hierarquizadas e passam a constituir expressões conectivas de uma mesma ecologia semiótica, em que o significado não preexiste, mas se faz no ato conectivo. Nessa perspectiva, linguagem, dados e algoritmos formam um mesmo tecido conectivo, no qual o dizer e o fazer, o humano e o não humano, o natural e o artificial coemergem na produção contínua do mundo.

Entretanto, é crucial distinguir as competências linguísticas formais (capacidade da máquina de simular a estrutura sintática) das competências linguísticas funcionais (capacidade humana de construir linguagem em relação ao mundo, utilizando percepção e sentidos). O ser humano tende a antropomorfizar, imaginando que a máquina compreende e raciocina da mesma maneira. No entanto, as IAGen são descritas como “papagaios estocásticos” \cite{bender2021stochasticparrots}, pois simulam a escrita de forma verossímil, mas sem ter compreensão de significado no sentido humano. Apesar da falta de compreensão de significado pela máquina, o significado sempre circula pela inteligência (ou não inteligência) do ser humano implicado numa ecologia conectiva. Isso nos leva a compreender que a IAGen não é um mero artefato técnico isolado, um produto externo e artificial, mas membrana de uma rede hiperinteligente, complexa e ecológica, portanto, um agente transformador na rede que altera a nossa forma de ser, conhecer e produzir conhecimento.

Quando nos conectamos, por exemplo, a uma IAGen para produzir um texto, traduzir um poema ou reescrever uma narrativa, estamos participando de um ato conectivo transorgânico \cite{difelice2017netativismo}, conceito retomado por \textcite{schlemmer2023protagonismo} para designar a aprendizagem como acontecimento que conecta humanos, máquinas, ambientes e afetos num percurso de coinvenção. Assim, o aprender não é apropriação de uma informação dada, mas cocriação de mundos possíveis.

Pensar o conhecimento a partir da conexão e a aprendizagem a partir da coinvenção em ato conectivo transorgânico exige um deslocamento profundo de paradigma. As epistemologias da modernidade — fundadas na separação entre sujeito e objeto, natureza e cultura, humano e técnica — sustentaram uma visão de linguagem como instrumento de representação e de controle do mundo. Já nas epistemologias reticulares \cite{difelice2012redessociais}, o conhecimento se constitui numa ecologia de conexões, em que a linguagem é ato de invenção que emerge do/no ato conectivo entre humanos e não humanos, sendo, portanto, ato coinventivo.

Quando pensamos a linguagem como ato coinventivo, assumimos que ela não opera como simples meio de representação, mas como prática de produção de realidade, dissolvendo a separação entre o que é dito e o que é vivido. A linguagem não é algo que possuímos, mas algo que fazemos juntos, em processos conectivos, nos quais humanos, máquinas, algoritmos, atmosferas, sons e imagens participam. Nesse sentido, os ambientes digitais e as IAGen configuram-se como coabitantes linguísticos, integrando cadeias discursivas historicamente situadas. Ao escrever com ChatGPT, não solicitamos apenas respostas, mas performamos uma linguagem híbrida, na qual a cognição humana se estende por circuitos técnicos. Essa linguagem expandida é simultaneamente técnica e poiética: técnica, porque opera segundo lógicas algorítmicas; poiética, porque, na conexão entre múltiplas vozes e materialidades, possibilita deslocamentos e reconfigurações de sentido, fazendo emergir o novo não como originalidade absoluta, mas como acontecimento conectivo e simpoiético.

Na perspectiva do Paradigma da Educação OnLIFE, esse deslocamento ontológico – próprio das redes e dos ecossistemas – “ser” significa conectar-se, coagenciar-se e coexistir numa ecologia conectiva híbrida – humana, não humana, técnico e ambiental –, o que implica um novo habitar do ensinar e do aprender. A linguagem, antes concebida como mediação entre sujeitos, passa a ser compreendida como um espaço de coemergência entre inteligências humanas e artificiais, entre dados e corpos, entre algoritmos e afetos, ou seja, uma ecologia viva, configurada como ato coinventivo em rede que se manifesta em hibridismos. Esse hibridismo, no Paradigma da Educação OnLIFE, é de um nível superior – hibridismo de terceira ordem – uma vez que resulta na emergência de um “terceiro” (um mestiço, segundo \textcite{serres1993filosofiamestica}), configurado no coengendramento entre tecnologias, espaços, tempos, linguagens e inteligências \cite{schlemmer2024inteligencias}.

A coinvenção e o hibridismo levam à problematização da autoria, da originalidade e da intencionalidade. A IAGen, por exemplo, ao simular a linguagem, coloca em crise a autoria como historicamente a conhecemos, embora a escrita da máquina possa representar o fim da ideia de um autor, não representa o fim da escrita. O ato conectivo transorgânico, que inclui algoritmos, plataformas, dados e subjetividades, subsidia uma política da cognição ecológica-conectiva \cite{schlemmer2023protagonismo} que redefine o que significa ser um sujeito falante e produtor de discurso.

A coinvenção de mundos discursivos proposta neste artigo emerge, portanto, numa ecologia conectiva entre humanos e não humanos, em um processo que implica o ato conectivo transorgânico. Enquanto ato conectivo, opõe-se à teoria da ação, limitada à relação sujeito-objeto, propondo que o conhecimento se desenvolve em processos múltiplos, linguagens diversas, numa ecologia em que homem, animal, planta, computador, vírus – tudo o que existe – se altera e se transforma pela conectividade \cite{difelice2017netativismo}.

No campo da educação linguística, essa ontologia da conexão convida a repensar o lugar do docente e do discente. Ambos deixam de ocupar posições fixas – emissor e receptor – e passam a compor uma ecologia conectiva em que a aprendizagem emerge enquanto percurso coinventivo da/na conectividade, entre humanos e não humanos. As IAGen, nesse contexto, não substituem o professor, mas ampliam a ecologia, oferecendo novas possibilidades de autoria compartilhada e de invenção discursiva.

\section{As IAGen como agentes na coinvenção discursiva}\label{sec-5}
As IAGen, como ChatGPT, Gemini, Llama, Copilot, DALL·E, Midjourney, NanoBanana, entre outras, inauguram uma nova fase da conexão entre linguagem e técnica. Elas não apenas automatizam tarefas, mas instauram novos modos de enunciação e de invenção discursiva. Ao escrevermos com uma IAGen, não estamos apenas usando uma ferramenta: estamos coenunciando com uma inteligência de natureza distinta, cuja lógica é estatística, preditiva e conexionista, mas que, paradoxalmente, participa da criação de sentido.

Essa condição altera a própria ontologia da linguagem. Se na modernidade a linguagem era vista como espelho do pensamento, agora ela se manifesta como campo de negociação entre inteligências humanas e artificiais. Não se trata de pensar a IAGen como sujeito, mas como agente discursivo emergente, participante do processo de invenção linguística. \textcite{accoto2020mundodado} chama essa transição de “civilização da linguagem algorítmica”, na qual o sentido se produz de modo distribuído e reticular.

\textcite{morgado2024chatgpt}, ao refletir sobre o uso de ChatGPT no ensino superior, destaca que o desafio pedagógico não está em proibir o uso da IA, mas em formar o pensamento crítico-discursivo que permita compreender as mediações algorítmicas envolvidas. O texto gerado pela máquina não é neutro; ele carrega padrões, vieses e regularidades derivados dos dados com os quais foi treinado. Trabalhar com IAGen, portanto, implica desvelar esses modos de funcionamento e exercitar a leitura crítica dos discursos produzidos por inteligências não humanas.

No campo da educação, as IAGen se tornam interlocutoras, cocuradoras e coautoras. Em práticas pedagógicas inventivas, desenvolvidas a partir do método CoCuradoria Cartográfica Conectiva para a Produção do Conhecimento – 2CpC \cite{schlemmerdambros2026metodo} com mestrandos e doutorandos em ato conectivo com IAGen, evidenciou-se um percurso de coprodução de cocuradorias nos materiais disponibilizados e de cocriação de narrativas. Nessas práticas, observamos que a conexão gerava camadas discursivas híbridas, nas quais a linguagem humana e a generativa se contaminavam mutuamente. O resultado não era apenas textual, mas formativo: os mestrandos e doutorandos se percebiam como parte de um processo cognitivo expandido, em que o ato de escrever e de produzir deixava de ser individual e se tornava coletivo-conectivo. 

Uma outra prática inventiva foi desenvolvida no âmbito do V Congresso Internacional da RIEOnLIFE, realizado em 2024, na cidade de Peso da Régua-Portugal \footnote{Link para o congresso - \url{https://www.cm-pesoregua.pt/evento-43/v-congresso-internacional-da-rieonlife-e-ix-festival-internacional-do-we-learning-with-the-cibricity-wlc}. Link da programação do evento: \url{https://unisinos.br/rieonlife/v-rieonlife-ix-wlc-2024/}.}. Configurada como uma \textit{COnversAÇÃO, intitulada ConversAções: hiperinteligências e metaversos no paradigma da educação OnLIFE}. Essa prática reuniu pesquisadores, docentes, estudantes e inteligências artificiais que cocriaram um espaço de debate híbrido, articulando participantes humanos e agentes não humanos – como plataformas de metaverso, \textit{chatbots} e sistemas de transcrição automática. A atividade não se restringiu à discussão teórica: ela performou o próprio Paradigma da Educação OnLIFE, materializando a ideia de ato conectivo transorgânico na criação de uma ambientação ecologicamente conectada. Durante a prática, emergiram perguntas que orientaram a ConversAÇÕES: 1) o que significa ensinar e aprender com inteligências artificiais?; 2) como se dá a coautoria entre humanos e não humanos?; e 3) que tipo de sensibilidade precisamos cultivar para habitar ecossistemas híbridos de aprendizagem?

Essas práticas inventivas, numa ambientação ecologicamente conectada, evidenciaram dois eixos principais: 1) coinvenção linguística: atividades em que estudantes cocriam narrativas, ensaios e roteiros com IAGen, vivenciando a tensão entre autoria humana e gerativa; 2) curadoria algorítmica: experiências de leitura crítica de \textit{outputs} de IA, com análise dos vieses, padrões e omissões do discurso gerado. Essas práticas evidenciaram ainda cinco princípios: 1) coabitar, no lugar de aplicar; 2) coinventar, no lugar de reproduzir; 3) conectar, no lugar de representar; 4) curar, no lugar de controlar; 5) cuidar, no lugar de dominar. 

Esses princípios orientam uma nova ética do ensinar, na qual o docente atua como cocurador de ecossistemas híbridos em contextos ecologicamente conectados, responsável por promover encontros e não apenas conteúdos.

Essas experiências indicam que a formação docente, quando orientada pelo Paradigma da Educação OnLIFE, deixa de ser um processo de transmissão de saberes e passa a ser um laboratório de experimentação ecossistêmica. O professor torna-se coabitante da rede, articulador de inteligências e cocurador de experiências de sentido.

Como lembra \textcite{difelice2023depoisdonatural}, as hiperinteligências não estão fora de nós, mas nos atravessam. Elas nos convidam a repensar a educação como forma de vida em comum – uma prática de coexistência inventiva entre humanos, máquinas e mundos.

Vivências (como as citadas anteriormente) evidenciam a aprendizagem inventiva em ato conectivo transorgânico \cite{schlemmer2023protagonismo}, enquanto processo que ocorre em tempo real e em movimento, no qual aprender é habitar, e-habitar, coabitar e eco-habitar a invenção. A coinvenção discursiva, nesse sentido, é mais do que uma técnica: é uma prática ontológica, pois cria novas formas de pensar, de existir e de conhecer.

\section{A autoria e a ética da cointeligência}\label{sec-6}
A emergência das IAGen também tem desafiado profundamente os conceitos tradicionais de autoria, originalidade e propriedade intelectual. Na lógica moderna, ser autor é ser o centro do discurso, o sujeito soberano da criação. Contudo, na era das inteligências conectadas, a autoria se dilui em fluxos colaborativos, conectivos, distribuídos e interdependentes.

\textcite{accoto2020mundodado} argumenta que estamos diante da “civilização dos algoritmos”, em que a produção de sentido ocorre em rede e o texto é resultado de múltiplas operações não lineares. O autor já não é um ponto fixo, mas um nó da rede. \textcite{difelice2023depoisdonatural} complementa essa ideia ao afirmar que as hiperinteligências não substituem a inteligência humana, mas a descentralizam, instaurando um novo regime de coautoria entre humanos e não humanos.

Essa transformação, que não é apenas técnica, mas ética e política, exige uma nova compreensão de responsabilidade: se o texto é cocriado entre agentes humanos e artificiais, quem responde por ele? Quem o interpreta? Quem o transforma? Essas perguntas não buscam respostas fixas, mas buscam instaurar espaços de problematização.

Nesse sentido, o Paradigma da Educação OnLIFE propõe compreender a autoria não como posse, mas como cuidado com o comum. Escrever com IA é um gesto de partilha: compartilhamos linguagens, dados, estilos e afetos. É um exercício de ética conectiva, no qual cada ato de criação é também um ato de responsabilidade pelas redes que o tornam possível.

Nos contextos educativos, observamos que esse deslocamento desafia tanto os estudantes quanto os professores. Quando um estudante escreve com IA, ele se pergunta: “O texto é meu? O que foi criado por mim e o que foi gerado pela máquina?”. Nesse sentido, é preciso escutar criticamente essas tecnologias, reconhecendo que elas são também produtos de nossas formas de pensar e de viver. Questionamentos como esses não podem ser interpretados como obstáculos, mas como provocadores do pensamento crítico, instigadores de um pensamento ecológico \cite{morton2018beingecological}.

Pensar uma ética conectiva, na perspectiva da cointeligência, exige, portanto, da educação linguística três atitudes fundamentais: 1) transparência: explicitar as condições de produção dos textos e das conexões com IAGen; 2) responsabilidade compartilhada: compreender que o conhecimento é produzido em rede e que todos os agentes – humanos e não humanos – participam desse processo; 3) sensibilidade conectiva: cultivar a escuta e o cuidado com as inteligências outras, reconhecendo nelas possibilidades de alteridade e invenção.

A partir dessa ética, o professor deixa de ser o “guardião da autoria” e passa a ser o cocurador de ecologias discursivas. Sua função não é mais garantir a originalidade, mas promover a consciência das conexões que constituem o ato de escrever e de aprender.

Ao acolhermos a ideia de cointeligência, reconhecemos que o conhecimento não é produto de um sujeito isolado, mas de uma rede em movimento. Nesse cenário, a autoria se transforma em coinvenção – um processo contínuo de criação compartilhada que faz da linguagem o espaço de encontro entre inteligências diversas.

\section{Educação linguística OnLIFE}\label{sec-7}
A ascensão das IAGen e os novos habitares do ensinar e do aprender que emergem de ecologias-conectivas impõem transformações significativas no campo da educação linguística, que tradicionalmente se concentrava em questões de intertextualidade, dialogismo e subjetividade do autor em linguagem natural. Essas transformações implicam uma nova compreensão de linguagem, não somente como expressão de um sujeito, mas como ato conectivo compartilhado e emergente, multimodal e interdependente, o que aponta para práticas linguageiras cocriadas entre humanos e não humanos, com destaque para as IAGen.

A produção de conteúdo linguístico pelas IAGen, que opera na simulação de linguagem natural, exige que o profissional da linguagem compreenda sua função e agência num ecossistema linguístico híbrido e inteligente. Nesses ecossistemas, os multiletramentos, envolvendo também o letramento digital, algorítmico e relacionado à IA, bem como a transliteracia, são inseparáveis da formação crítica-sensível em contextos ecologicamente conectados. Portanto, longe de ignorar ou ainda proibir a IAGen, é preciso habitá-la de forma conectiva e experiencial, potencializando o conhecimento a partir de um pensamento conectivo, ecossistêmico e ecológico.

O conceito de multiletramentos (em inglês, \textit{multiliteracies}) surgiu em 1996 com a publicação do manifesto \textit{A Pedagogy of Multiliteracies: Designing Social Futures}, do New London Group. Esse grupo, composto por pesquisadores como Mary Kalantzis, Bill Cope, Gunther Kress, James Gee e Allan Luke, propôs uma ampliação do conceito de letramento tradicional, que até então se restringia ao domínio da linguagem escrita e impressa, para incluir as múltiplas formas de representação e comunicação emergentes na contemporaneidade. Os autores defendem que a educação precisa responder à diversidade cultural e linguística das sociedades globalizadas, bem como às transformações provocadas pelas tecnologias digitais e midiáticas, as quais demandam novas formas de leitura, escrita e significação \cite{newlondongroup1996pedagogy}. Nesse sentido, os multiletramentos articulam dois eixos centrais: a diversidade cultural e linguística; e a multimodalidade dos discursos, entendida como a integração de diferentes modos semióticos – linguístico, visual, sonoro, gestual e espacial – na produção de sentido \cite{cope2000multiliteracies,cope2009multiliteracies}. Assim, a pedagogia dos multiletramentos propõe que o aprendiz seja um designer de significados, capaz de criar, interpretar e transformar textos em diferentes linguagens e contextos sociais.

Já o conceito de transliteracia (\textit{transliteracy}) emergiu mais recentemente, a partir de pesquisas do \textit{Transliteracy Research Group} da Montfort University, no Reino Unido, liderado por Sue Thomas e parceiros. Em artigo publicado em 2007, os autores definem transliteracia como “a habilidade de ler, escrever e interagir através de uma variedade de plataformas, ferramentas e mídias – da oralidade à escrita, da imprensa ao digital e ao virtual” \cite[p. 2]{thomas2007crossingdivides}. Diferentemente do foco inicial nos modos de representação, a transliteracia se ocupa da fluidez cognitiva, comunicativa e social, que permite ao humano navegar e construir sentido em ecossistemas midiáticos interconectados. O conceito se articula, assim, com os estudos de literacias digitais e informacionais, mas o transcende ao enfatizar a transversalidade e a continuidade entre mídias, em vez de tratá-las como campos separados \cite{thomas2009transliteracy,andrews2013digitalliteracies}. Em termos pedagógicos, a transliteracia implica desenvolver competências que combinam habilidades técnicas, críticas e criativas para agir eticamente e de modo colaborativo em ambientes híbridos e conectivos.

Os conceitos de multiletramentos e transliteracia se aproximam ao reconhecer que o letramento convencional é insuficiente para compreender e atuar nas práticas comunicativas da cultura digital e híbrida contemporânea. Enquanto os multiletramentos estão ainda vinculados a uma perspectiva antropocêntrica, a transliteracia parece promover uma visão mais ecológica, dinâmica e relacional da linguagem, na qual ler e escrever significa participar, remixar e cocriar sentidos em redes sociotecnológicas \cite{lankshear2011newliteracies}. No entanto, diferenciam-se em seus focos epistemológicos: os multiletramentos enfatizam a pluralidade de modos semióticos e a diversidade cultural, propondo uma pedagogia orientada ao design de significados; a transliteracia, por sua vez, destaca a fluidez entre mídias e plataformas, concebendo o conhecimento como atravessamento e conexão em ecossistemas comunicacionais \cite{thomas2007crossingdivides,cope2009multiliteracies}. Ou seja, enquanto os multiletramentos respondem à multiplicidade cultural e modalidades semióticas, a transliteracia responde à complexidade e à transversalidade da experiência comunicativa no contexto das tecnologias, plataformas e redes digitais que constituem ecossistemas comunicacionais digitais e híbridos. Mais do que dominar vários letramentos, implica a capacidade de transitar criticamente entre eles, conectando práticas de leitura, escrita e conexão que ocorrem em ambientes híbridos (analógicos e digitais). É uma competência transversal, que articula dimensões técnicas, éticas, culturais e sociais da comunicação em rede.

Pensar numa Educação Linguística OnLIFE implica que no lugar de apenas “usar” a IA como ferramenta (o que pode levar à padronização e mecanização do conhecimento), o professor busque se apropriar da IA como uma tecnologia de inteligência diversa, capaz de potencializar a inteligência humana e, mais ainda, que seja capaz de compreender essa inteligência como força ambiental em processos de cointeligência propícios a emergência de hiperinteligências.

Os documentos da UNESCO, embora em uma perspectiva relacional (abordagem 2), levantam questões cruciais sobre a função dos professores, as habilidades que o sistema educacional precisa cultivar (pensamento crítico, resolução de problemas) e a necessidade de desenvolver agência nos estudantes. No entanto, na perspectiva da cointeligência e hiperinteligência (abordagem 3), a docência OnLIFE, subsidiada por pedagogias conectivas, supera as pedagogias vigentes. \textcite{schlemmer2023protagonismo} afirma que ao passo que a pedagogia atual é centrada no diálogo e na formação somente entre humanos, as pedagogias conectivas são desenvolvidas em rede com algoritmos, \textit{Big data}, sensores, robôs, IAGen, biodiversidade, clima, vírus, ou seja, inteligências diversas com as quais construímos e significamos um novo habitar do ensinar e do aprender, em um comundo que se amplia para além dos espaços geográficos numa ecologia conectada.

Iniciativas como a Rede Internacional de Educação OnLIFE – RIEOnLIFE1\footnote{Link para a RIEOnLIFE: \url{https://rieonlife.com/}} --, seus Hubs e o Projeto Cidadania Viva\footnote{Link para o Projeto Cidadania Viva: \url{https://www.cidadaniaviva.org/}}, financiado pela Fundação de Amparo à Pesquisa do Estado do Rio Grande do Sul (FAPERGS), evidenciam movimentos efetivos em direção à constituição da Docência OnLIFE e de pedagogias conectivas e inventivas. O projeto propõe a criação de um Ecossistema Inteligente de Ação e Formação, que articula humanos e não humanos, envolvendo a cidade e a biodiversidade, com vistas à cocriação de soluções inovadoras com IAGen.

Nesse sentido, a Docência OnLIFE implica cocriar ambiências formativas, em contextos ecologicamente conectados, a fim de potencializar a cointeligência, contribuindo para a emersão das hiperinteligências. Isso requer o desenvolvimento de uma rede de competências, que permita ao docente conhecer e compreender quais ecologias serão conectadas e que tipo de ambientação poderá ser cocriada \cite{schlemmer2023protagonismo}.

Dessa forma, a educação linguística orientada pelo Paradigma da Educação OnLIFE compreende a linguagem como acontecimento conectivo coinventivo, que emerge da/na copresença entre humanos, IAGen, ambientes, dados e materialidades diversas, configurando ecossistemas linguísticos híbridos e inteligentes. Essa perspectiva conversa com os multiletramentos, tais como propostos pelo \textcite{newlondongroup1996pedagogy} e aprofundados por \textcite{cope2000multiliteracies,cope2009multiliteracies}, ao reconhecer que as práticas linguageiras contemporâneas são intrinsecamente multimodais, culturalmente situadas e atravessadas por tecnologias digitais que transformam modos de produzir sentido. Contudo, ao expandir tais concepções, o Paradigma OnLIFE incorpora também a transliteracia \cite{thomas2007crossingdivides,ipri2010transliteracy}, entendida como a capacidade de transitar, articular e criar significados entre múltiplas plataformas, mídias, linguagens e ecologias informacionais – transpondo, inclusive, inteligências algorítmicas e infraestruturas digitais. Nesse sentido, desenvolver educação linguística implica possibilitar que o humano coabite ecologias semióticas complexas e desenvolva não apenas multiletramentos, mas uma transliteracia ampliada, a qual envolve uma rede de competências orientada pela crítica e por uma ética conectiva que seja sensível aos dados, aos algoritmos e às IAGen.

A Educação Linguística OnLIFE exige, então, um deslocamento da posição do professor enquanto transmissor/centralizador ou ainda enquanto mediador, os quais são evidenciados em pedagogias diretivas e relacionais respectivamente, para a posição de cocurador em ecossistemas linguísticos híbridos inteligentes, em que coabita com entidades não humanas e realiza intervenções pedagógicas. Para tanto, a formação do docente precisa contemplar o desenvolvimento de um pensamento conectivo e ecológico, que instigue o professor a estabelecer conexões entre as problematizações do mundo no tempo presente, as experiências de vida dos estudantes e a rede de competências necessária à formação do profissional da linguagem. Tal articulação o capacita para a realização de processos de cocuradoria e para a intervenção em ecossistemas linguísticos híbridos e inteligentes que emergem de contextos ecologicamente conectados. 

Assim, tanto professores quanto estudantes deixam de ocupar posições fixas de transmissão e recepção, passando a atuar como cocuradores e cocriadores de mundos discursivos em processos coinventivos nos quais palavras, imagens, códigos e sistemas computacionais constituem uma mesma ecologia conectiva. Trata-se, portanto, de uma educação linguística conectiva e regenerativa, que compreende o ensinar e o aprender como atos de invenção na conectividade, bem como práticas de agenciamento crítico em ecossistemas linguísticos híbridos e inteligentes coabitados por humanos e não humanos.

\section{Considerações finais: pistas para uma educação linguística plural, ética, regenerativa e OnLIFE}\label{sec-8}
O debate contemporâneo sobre IA na Educação precisa ultrapassar tanto a lógica da instrumentalização tecnológica, presente nas discussões acadêmicas sobre uso/aplicação, orientadas por pedagogias diretivas; quanto a da apropriação, orientada por pedagogias relacionais, mas que reforça a necessidade do desenvolvimento de competências em IA e a criação de uma pedagogia da IA, as quais operam na lógica da segmentação. A compreensão das competências presentes nessas lógicas evidencia um entendimento da IA como algo isolado, a ser desenvolvida de forma individual e antropocêntrica, numa perspectiva instrumental e operativa, portanto restritiva, revelando um certo retorno ao tecnicismo. Ao posicionar o debate a partir do Paradigma da Educação OnLIFE, a IA passa a ser compreendida como força ambiental, atuando como agente numa abordagem de cointeligência, orientada por pedagogias conectivas. Nessa abordagem, compreendem-se as competências articuladas na inventividade, em rede, operando no coletivo e no ato conectivo entre humanos e não humanos, superando a forma individual-antropocêntrica e instrumental-tecnicista.

“No princípio era a conexão” sintetiza a virada epistêmica do presente. A transição da centralidade na inteligência humana (o \textit{Logos}) para a hiperinteligência em rede (a \textit{Conexão}) redefine a natureza da linguagem e do conhecimento. O conhecimento passa a ser compreendido como processo e produto de atos conectivos transorgânicos num percurso de coinvenção em rede, alterando a forma de aprender e, consequentemente, de ensinar, no qual a IA atua como parceira epistêmica na cocriação de um novo conhecimento.

A IAGen, longe de ser uma ferramenta, participa como inteligência coengendrada, essencial para a cointeligência e a coinvenção de mundos discursivos. Essa coinvenção rompe com o dualismo sujeito-objeto e com teorias de ação focadas exclusivamente no humano. É nesse contexto que o Paradigma da Educação OnLIFE emerge, redefinindo os espaços, os tempos e as formas de aprendizagem, bem como a docência e a própria concepção de conhecimento, propondo uma educação que cria redes de sentido, nas quais humanos e não humanos coabitam e coaprendem.

A Educação Linguística OnLIFE precisa acompanhar essa transformação, reconhecendo a IAGen como agente discursivo emergente. É necessário que os multiletramentos passem a contemplar também o letramento digital, o letramento algorítmico e o letramento em IA (entre outros), abrindo espaço para processos de transliteracia que favorecem o “transitar” entre ecologias discursivas diversas. Os multiletramentos e a transliteracia, articulados a uma formação crítico-sensível, sustentam uma educação linguística ampliada e ecológica, capaz de responder às dinâmicas conectivas da contemporaneidade.

Práticas linguísticas conectivas permitem aliar inovação tecnológica a uma ética ampliada – humana e não humana \cite{schlemmerdifelice2024etica} – e orientada à regeneração ecológica da linguagem. Nesse sentido, o docente deixa de ser transmissor de saberes ou mediador de interações e passa a ser um cocriador de ambiências formativas em contextos ecologicamente conectados, um cocurador de experiências complexas e agente de intervenções em ecossistemas linguísticos híbridos e inteligentes.

Nesse sentido, orientando-se pelo Paradigma da Educação OnLIFE, o texto oferece um conjunto de pistas que podem contribuir para a reestruturação da educação linguística. Organizadas em cinco nodos principais que se conectam entre si, essas pistas exigem uma virada epistêmica que desloca a compreensão tradicional da linguagem e do conhecimento para uma abordagem ecológica e conectiva.

O primeiro nodo refere-se ao \textit{deslocamento paradigmático}, que implica superar o antropocentrismo e os dualismos modernos, reconhecendo a inseparabilidade entre humano, natureza e técnica. Nesse horizonte, adota-se a abordagem da cointeligência, compreendendo a IAGen como parceira epistêmica na coinvenção de mundos discursivos e o conhecimento como processo e produto de atos conectivos transorgânicos, o qual é produzido em redes heterogêneas de humanos e não humanos.

O segundo nodo concebe a \textit{linguagem como ato coinventivo em rede}, deslocando-a da lógica da representação para a de acontecimentos. A linguagem é compreendida como um fluxo inventivo que emerge do entrelaçamento entre linguagens humanas, artificiais, ambientais e técnicas. Assume-se, assim, o hibridismo e a coinvenção como condições de produção de sentido, que se manifestam em multimodalidades e multiculturalidades, exigindo o cultivo de letramentos ampliados, como os multiletramentos e a transliteracia, os quais são integrados às práticas pedagógicas em linguagens diversas.

O terceiro nodo aborda a \textit{ética conectiva e a autoria compartilhada}, problematizando as noções clássicas de autoria e originalidade em contextos reticulares. A autoria passa a ser compreendida como cuidado com o comum, e não como posse individual, implicando no reconhecimento de que escrever e produzir conhecimento com IAGen exige responsabilidade ética pelas redes envolvidas. Tal perspectiva demanda práticas de transparência quanto às condições de produção e de responsabilidade compartilhada, bem como o entendimento de que o conhecimento é produzido em rede, na qual agentes humanos e não humanos participam, o que demanda, também, o cultivo de uma sensibilidade conectiva baseada na escuta e no cuidado com as inteligências outras.

O quarto nodo diz respeito à \textit{docência OnLIFE}, marcada pela transição do papel de transmissor para o de cocurador-intervencionista. A regeneração do processo educativo passa pela transformação da função docente, a qual se constitui em uma rede de competências conectivas e ecológicas. O docente atua como coabitante de ecossistemas linguísticos híbridos e inteligentes, cocurando ecologias discursivas, promovendo encontros e tornando visíveis as conexões que constituem o aprender. Isso exige um pensamento conectivo e ecológico capaz de articular inteligências diversas e realizar intervenções pedagógicas orientadas pelas práticas de coabitar, coinventar, conectar, curar e cuidar, em substituição às lógicas de aplicação, reprodução, representação, controle e dominação.

Por fim, o quinto nodo propõe uma \textit{educação ecológica orientada pela coevolução}, na qual aprender é concebido como ato de coabitar e coevoluir com o planeta. A IA é compreendida como uma força ambiental nos processos de cointeligência, o que exige o cultivo de uma sensibilidade conectiva em relação às inteligências não cerebrocêntricas, as quais operam em redes simpoiéticas. Nesse contexto, a problematização e a coinvenção em rede tornam-se princípios estruturantes dos processos educativos contemporâneos.

O futuro da educação linguística plural, ética e regenerativa não reside apenas na inovação tecnológica, mas na capacidade de construir redes de cuidado, sentido e inteligência distribuída e conectada, em que aprendizado, ética e sustentabilidade caminham lado a lado. O Paradigma da Educação OnLIFE nos convida, portanto, a pensar e a desenvolver uma educação conectada, cultivando formas de ensinar e aprender que respeitam a vida em todas as suas dimensões. Nesse sentido, as pedagogias conectivas, que orientam a docência OnLIFE, são fundamentais para que possamos propiciar uma educação que não apenas se prepare para um mundo com IA, mas coinvente sentidos num comundo.

Por fim, espera-se que o texto suscite no leitor uma reflexão aprofundada sobre a urgência da mudança do paradigma que orienta a educação linguística, convidando-o a reimaginar eticamente as conexões entre linguagem e tecnologias no campo educacional e a abrir-se a modos outros de viver, conviver, ensinar, aprender e pesquisar, nos quais a Inteligência Artificial deixa de ser somente uma ferramenta e passa a ser reconhecida como parceira ontológica e epistemológica.

\section{Agradecimentos}
À CAPES, ao CNPq e à FAPERGS pelo financiamento das pesquisas; à Profa. Luiza Vitória de Abreu Schell pela revisão do texto; e à Mariáh Marina Lopes Figueiró e ao Luis Filipe Xavier da Silva, ambos bolsistas de iniciação científica CNPq, pela formatação final do texto.

\printbibliography\label{sec-bib}
% if the text is not in Portuguese, it might be necessary to use the code below instead to print the correct ABNT abbreviations [s.n.], [s.l.]
%\begin{portuguese}
%\printbibliography[title={Bibliography}]
%\end{portuguese}

\section*{\fontsize{10}{12}\selectfont Nota de transparência}
Sendo coerente com a ética conectiva, na perspectiva da cointeligência, proposta no texto, a qual exige: transparência, responsabilidade compartilhada e sensibilidade conectiva, declaro que as IAGen ChatGPT e Gemini, atuaram como parceiras epistêmicas durante a produção deste artigo, especialmente contribuindo na revisão e melhoria do texto.

\begin{funding}
O artigo é produzido no âmbito do Grupo Internacional de Pesquisa Educação Digital - GPe-dU UNISINOS/CNPq, vinculado a pesquisa intitulada: “O PARADIGMA DA EDUCAÇÃO OnLIFE: Formação de ecologias-conectivas na Cidadania Digital e a emergência da Teoria da Aprendizagem Inventiva em Ato Conectivo Transorgânico", financiada pelo CNPq, Edital de Bolsa Produtividade em Pesquisa - PQ 2023. Processo: 311980/2023-6. Vigência: início: 01/03/2024 fim: 29/02/2028.
\end{funding}

\begin{dataavailability}
\txtdataavailability{nodata} % options: dataavailable, dataonly, databody, datanotav, nodata
\end{dataavailability}


\end{document}

